\documentclass[11pt,a4paper]{article}
\input{T0_preamble_local_De}
\hfuzz=8pt\tolerance=5000\emergencystretch=3em
\title{Algebraische Brücken zur Bewusstseinsdiskussion\\
\large Anschluss an die Dokumente 100 und 170}
\author{Johann Pascher\\
\small Independent Researcher, Oberösterreich, Austria\\
\small ORCID: \href{https://orcid.org/0009-0000-6518-4064}{0009-0000-6518-4064}\\
\small \href{mailto:johann.pascher@gmail.com}{johann.pascher@gmail.com}}
\date{8. Mai 2026}
\begin{document}
\maketitle
\begin{abstract}
\noindent
Dokument 206 (\emph{Dreieck-Matrix-Reduktionstheorem}) liefert eine algebraische Struktur, die die bisherige Bewusstseinsdiskussion in der FFGFT-Reihe geometrisch erdet, ohne sie inhaltlich zu erweitern. Die phänomenologische Linie ist in Dokument~100 (\emph{Bewusstsein als fraktale Inkohärenz}, Anschluss an das No-Go-Theorem von Adlam, McQueen und Waegell, 2025) und in Dokument~170 (\emph{Bewusstsein als topologisch erzwungene Schwellenleiter}) bereits gezogen.

Wir zeigen drei strukturelle Brücken: (A) Tekermens \emph{konstitutive Offenheit} ist algebraisch identisch mit dem Nichtabschluss-Theorem (Doc~193) und entspricht der Bedingung $\xi > 0$; ein perfekt unitäres System hat $\det A = 1$, also keine Welt-Beziehung. (B) Phillips' \emph{Information Grounding Law} ($A^2 = A^\top$) ist die Z$_3$-Selbstinversion und entspricht der \emph{rekursiven sensorischen Rückkopplung} aus Dokument~170. (C) Die Z$_3^3$-Skalenleiter mit Eigenwerten $(1-\xi)^{n/3}$ liefert das algebraische Skelett für die diskreten Komplexitätsstufen, die Dokument~170 als Arbeitshypothese formuliert; sie erzwingt Diskretheit, aber nicht die phänomenologische Stufenzuordnung.

Wir grenzen scharf ab, was die Algebra \emph{nicht} leistet: keine Ableitung von Bewusstsein aus $\xi$, keine Vorhersage, welche Schicht $n/3$ bewusst ist, keine Auflösung des Hard~Problem. Die Brücken sind strukturell, nicht ontologisch. Ein methodischer Abschnitt überträgt die Selbstkontrollen aus Dokument~206 (Numerologie-Filter, Zirkularitätsabschnitt) auf das Bewusstseinsthema. Ein zusätzliches Kapitel formuliert vier qualitative empirische Vorhersagen unter der Diskretheits-Hypothese aus Dokument~170 --- Bimodalität, Entwicklungssprung, Anästhesie-Kippen, Spezies-Klassifikation --- jeweils ohne quantitative Zahlenwerte aus $\xi$.

Die Bilanz: Dokument~206 macht die Sprache der Diskussion algebraisch konsistent, aber die Bewusstseinsfrage selbst bleibt eine offene Forschungsfrage; sie wird durch Doc~206 weder gelöst noch geschlossen.
\end{abstract}
\tableofcontents
\newpage
% ============================================================
% dok207_bewusstsein_bruecken_De_ch.tex
% Algebraische Brücken zur Bewusstseinsdiskussion
% FFGFT Dokument 207 — Inhaltsteil
% Status: 8. Mai 2026
% ============================================================

\section{Einführung}\label{sec:einleitung}

Die Bewusstseinsdiskussion ist innerhalb der FFGFT-Reihe seit längerem etabliert. Zwei Dokumente bilden ihre Hauptlinien:

\begin{itemize}
\item \textbf{Dokument 100 (\emph{Bewusstsein als fraktale Inkohärenz}):} Anschluss an das No-Go-Theorem von Adlam, McQueen und Waegell~\cite{Adlam2025}, das zeigt, dass Agentität in einem rein unitären Quantensystem nicht entstehen kann (Weltmodell-Versagen, Deliberations-Unmöglichkeit, Aktionsselektion-Zusammenbruch). FFGFT antwortet: Bewusstsein entsteht nicht aus perfekter Quantenkohärenz, sondern aus \emph{strukturierter fraktaler Inkohärenz} — einer geometrisch verwurzelten Abweichung von Unitarität, parametrisiert durch $\xi = 4/30{,}000$.
\item \textbf{Dokument 170 (\emph{Bewusstsein als Schwellenleiter}):} Bewusstsein als Stufe~4 einer fünfstufigen Hierarchie diskreter Komplexitätsschwellen, getragen von der fraktalen Korrektur $K_{\text{frak}}^{(n)} = (\xi/(1+\xi))^{n/2}$ als Schwellenwertoperator. Selbsterklärt als Arbeitshypothese.
\end{itemize}

Dokument~206 (\emph{Dreieck-Matrix-Reduktionstheorem}, ebenfalls 8.\,Mai 2026) hat mit dem Beweis $\det(A_\xi) = 1-\xi$ und der Z$_3^3$-Skalenleiter $(1-\xi)^{n/3}$ eine algebraische Struktur etabliert, die die phänomenologischen Aussagen der Dokumente~100 und~170 strukturell erdet. Parallel dazu hat die Diskussion in der \emph{Information Physics Institute}-Mailingliste in den Wochen vor diesem Dokument zwei zentrale Begriffe geprägt, die direkt in dieselbe Struktur übersetzbar sind: Onur Tekermens \emph{konstitutive Offenheit} und Phillips Pauls \emph{Information Grounding Law} (IGL) im Rahmen der Self-Dual~C-25-Konstruktion.

\medskip

Die Aufgabe dieses Dokuments ist begrenzt: Wir zeigen, dass die Sprache, in der bisher über Bewusstsein in FFGFT und in der IPI-Diskussion gesprochen wurde, mit dem in Doc~206 etablierten algebraischen Skelett konsistent ist. Wir formulieren keine neue Bewusstseinstheorie. Wir öffnen keine neuen Forschungsfragen, die nicht bereits in Dok~100 oder Dok~170 stehen. Wir prüfen lediglich, was die Algebra strukturell trägt.

\begin{important}[Was dieses Dokument ist und was nicht]
\textbf{Es ist:} Eine Anschluss-Schrift, die zeigt, wo Doc~206 die phänomenologische Linie aus Doc~100 und Doc~170 algebraisch trifft.

\textbf{Es ist nicht:} Eine Ableitung von Bewusstsein aus $\xi$. Die Algebra erzwingt Strukturmerkmale (Diskretheit, Offenheit, Selbstinversion); sie identifiziert keine Stufe als \glqq bewusst\grqq{}.
\end{important}

\section{Was Dokument 206 algebraisch leistet (Zusammenfassung)}\label{sec:doc206_zusammenfassung}

Für die Anschlussfähigkeit fassen wir die für dieses Dokument relevanten Resultate aus Doc~206 kompakt zusammen.

\subsection{\texorpdfstring{Z$_3$-Dreieck und 3$\times$3-Matrix}{Z3-Dreieck und 3x3-Matrix}}

Der Z$_3$-Zyklus hat Adjazenzmatrix
\[
A = \begin{pmatrix} 0 & 0 & 1 \\ 1 & 0 & 0 \\ 0 & 1 & 0 \end{pmatrix},
\quad
A^3 = I, \quad \det(A) = 1, \quad \mathrm{Tr}(A) = 0.
\]
Eigenwerte sind die dritten Einheitswurzeln $\{1, \omega, \omega^2\}$.

\subsection{\texorpdfstring{$\xi$-Störung und gestörte Adjazenz}{xi-Störung und gestörte Adjazenz}}

Die Schließkante wird mit $\xi$ gedämpft:
\[
A_\xi = \begin{pmatrix} 0 & 0 & 1-\xi \\ 1 & 0 & 0 \\ 0 & 1 & 0 \end{pmatrix},
\quad
A_\xi^3 = (1-\xi)\,I, \quad \det(A_\xi) = 1-\xi.
\]
Die fraktale Dimension folgt in linearer Ordnung: $D_f = 3 - \xi$.

\subsection{\texorpdfstring{Z$_3^3$-Tensorprodukt und Skalenleiter}{Z3-hoch-3-Tensorprodukt und Skalenleiter}}

Das Tensorprodukt $T(\xi) = A_\xi \otimes A_\xi \otimes A_\xi$ hat $27$ Eigenwerte vom Modul $(1-\xi)$ und liefert die Skalenleiter
\[
\lambda_n = (1-\xi)^{n/3}, \quad n = 1, 2, 3, \dots
\]
Diese Drittelpotenzen sind nicht willkürlich, sondern Folge der dreifachen Tensorstruktur.

\subsection{Sechs Brücken zu fremden Formulierungen}

Doc~206 testet sechs Brücken (Austin, Tekermen, Phillips, Porter, Chekanov-Hakan, Moseley). Für die Bewusstseinsfrage relevant sind \textbf{Tekermen} (konstitutive Offenheit) und \textbf{Phillips} (IGL als Self-Dual-Bedingung). Die übrigen Brücken werden hier nicht vertieft.

\section{Brücke A: Konstitutive Offenheit und Welt-Beziehung}\label{sec:brueckeA}

\subsection{Tekermens Postulat in der IPI-Diskussion}

Onur Tekermen formuliert in seinem Unified Information Field-Rahmen (UIFT) die These, dass eine Rekursion nur dann \emph{informationsfähig} ist, wenn sie nicht schließt. Ein vollständig geschlossenes System hat keine Welt, kein Außen, keinen Bezug. Tekermen nennt dies \emph{konstitutive Offenheit}.

\subsection{Algebraische Identifikation in FFGFT}

In FFGFT entspricht das exakt dem Nichtabschluss-Theorem aus Doc~193, das in Doc~206 \S6 algebraisch verankert wurde:
\begin{equation}
\det(A_\xi) = 1 - \xi.
\end{equation}
Für $\xi = 0$ wäre $\det = 1$, das System wäre exakt zyklisch (perfekt unitär in der Matrix-Lesart). Erst $\xi > 0$ öffnet den Zyklus algebraisch.

\begin{insightBox}[Algebraische Form von Tekermens Postulat]
Konstitutive Offenheit $\equiv$ $\xi > 0$. Die Rekursion ist sinnvoll genau dann, wenn $\det(A_\xi) < 1$.
\end{insightBox}

\subsection{Anschluss an Dokument 100 und an Adlams No-Go-Theorem}

Adlam, McQueen und Waegell zeigen, dass ein rein unitäres Quantensystem keine Agentität tragen kann. Dokument~100 antwortet: Bewusstsein entsteht nicht aus Quantenkohärenz, sondern aus fraktaler Inkohärenz.

In der Sprache von Doc~206 ist das algebraisch lesbar:

\begin{itemize}
\item \textbf{Adlams No-Go:} Unitäres System $=$ geschlossene Rekursion $=$ $\det = 1$, also $\xi = 0$.
\item \textbf{FFGFT-Antwort:} $\xi > 0$ erzwingt geometrisch eine Welt-Beziehung. Das System hat ein Außen, weil sein algebraisches Skelett nicht schließt.
\end{itemize}

Das No-Go-Theorem von Adlam und das Nichtabschluss-Theorem von FFGFT sind nicht zwei Theoreme, sondern zwei Seiten derselben strukturellen Aussage. Adlam zeigt, was \emph{nicht} geht in einem geschlossenen System; FFGFT zeigt, was im offenen System der Fall ist.

\begin{interpretation}
Diese Identifikation löst das Bewusstseinsproblem nicht. Sie sagt nicht, dass $\xi > 0$ Bewusstsein \emph{erzeugt}. Sie sagt nur: die strukturelle Vorbedingung für Welt-Bezug, die Adlam fordert und Tekermen postuliert, ist in FFGFT algebraisch erfüllt. Die Frage, was darauf aufbaut, bleibt offen.
\end{interpretation}

\section{Brücke B: Self-Dual und rekursive Rückkopplung}\label{sec:brueckeB}

\subsection{Phillips' Information Grounding Law}

Phillips Paul postuliert in der Self-Dual~C-25-Konstruktion eine algebraische Bedingung für \emph{informationsverankerte} Strukturen: Eine Operation $A$ heißt \emph{self-dual}, wenn
\begin{equation}
A^2 = A^\top = A^{-1}.
\end{equation}
Phillips nennt dies das \emph{Information Grounding Law}: Information ist nur dann verankert, wenn die zugrundeliegende Operation auf sich selbst zurückgreifen kann, ohne Informationsverlust.

\subsection{Algebraische Identifikation in FFGFT}

In Doc~206 \S7 ist gezeigt:
\begin{equation}
A^2 = A^\top \quad \text{für die Z}_3\text{-Adjazenzmatrix},
\end{equation}
und allgemeiner gilt $A^k = A^{\top, k}$ für alle ganzen $k$. Die Z$_3$-Selbstinversion ist die elementare algebraische Verkörperung von Phillips' IGL.

\subsection{Anschluss an Dokument 170: Rekursive sensorische Rückkopplung}

Dokument~170 nennt als erstes Kriterium für die Komplexitätsschwelle Stufe~2 (zelluläre Selbsterhaltung) die \emph{rekursive sensorische Rückkopplung}: Membranrezeptoren, chemische Gradienten und intrazelluläre Signalkaskaden bilden einen geschlossenen Rückkopplungskreis. Die Zelle ertastet ihre Umwelt nicht passiv, sondern moduliert ihre eigene Reaktion auf Basis vergangener Zustände.

In der Sprache von Doc~206 ist das die operative Bedeutung von $A^2 = A^\top$: Die Operation \glqq Wahrnehmen\grqq{} ist invertierbar zu sich selbst — Wahrnehmen und Reagieren sind nicht zwei getrennte Prozesse, sondern dieselbe Z$_3$-Operation in unterschiedlicher Richtung.

\begin{insightBox}[Algebraische Form rekursiver Rückkopplung]
Rekursive sensorische Rückkopplung $\equiv$ Self-Dual ($A^2 = A^\top$). Phillips' IGL und Dok~170 \S{}Rekursive sensorische Rückkopplung beschreiben dieselbe algebraische Eigenschaft.
\end{insightBox}

\subsection{Was die Brücke leistet — und was nicht}

Sie leistet: Sie zeigt, dass die Phänomenologie aus Dok~170 und das Postulat aus der IPI-Diskussion auf demselben Z$_3$-Skelett stehen. Eine \glqq rekursive sensorische Rückkopplung\grqq{} ist nicht poetisch oder metaphorisch — sie hat eine konkrete algebraische Form.

Sie leistet \emph{nicht}: Sie sagt nicht, dass alle Z$_3$-Selbstinversionen Bewusstsein tragen. Eine rotierende Schraube erfüllt $A^2 = A^\top$ und ist nicht bewusst. Self-Dual ist eine \emph{notwendige} strukturelle Bedingung gemäß Doc~170 und Phillips IGL, keine hinreichende.

\section{\texorpdfstring{Brücke C: Z$_3^3$-Skalenleiter und diskrete Schwellen}{Brücke C: Z3-hoch-3-Skalenleiter und diskrete Schwellen}}\label{sec:brueckeC}

\subsection{Doc 170: Tabelle der Komplexitätsstufen}

Dokument~170 schlägt eine fünfstufige Hierarchie vor:

\begin{center}
\begin{tabular}{c|l|l}
\textbf{Stufe} & \textbf{System} & \textbf{Schwellenmerkmal} \\
\hline
0 & Elementarteilchen & $T = 1/m$ als passives Zeitfeld \\
1 & Atome / Kristalle & Diskrete geometrische Ordnung \\
2 & Zelle & Rekursive Selbsterhaltung, Metabolismus \\
3 & Nervensystem & Integrierte sensorische Rückkopplung \\
4 & Bewusstsein & Selbstmodellierung, Welt-Bezug \\
\end{tabular}
\end{center}

Doc~170 markiert dies explizit als Arbeitshypothese.

\subsection{Doc 206: Algebraische Skalenleiter}

Doc~206 \S4 zeigt, dass die Z$_3^3$-Tensorstruktur Eigenwerte
\[
\lambda_n = (1-\xi)^{n/3}, \quad n \in \mathbb{Z}_{\geq 0}
\]
hat. Diese Drittelpotenzen sind nicht willkürlich gewählt, sondern algebraisch erzwungen durch die dreifache Tensorstruktur.

\subsection{Was die Algebra für Doc 170 leistet}

\begin{itemize}
\item \textbf{Diskretheit:} Die Stufen in Doc~170 sind diskret. Das ist mit der Z$_3^3$-Skalenleiter konsistent: $(1-\xi)^{n/3}$ ist nur für $n = 0, 1, 2, 3, \dots$ algebraisch verankert. Kontinuierliche Zwischenwerte haben keinen geometrischen Status.
\item \textbf{Hierarchie:} Höhere $n$ entsprechen kleineren Modulen $(1-\xi)^{n/3}$, also tieferer Rekursion. Das ist mit der Vorstellung kompatibel, dass komplexere Strukturen tiefer eingebettet sind.
\item \textbf{Topologische Erzwingung:} Doc~170 sagt, die Stufen seien \glqq topologisch erzwungen\grqq{}. Das ist algebraisch verankert: die $\omega^k$-Phasen mit $k = 0, 1, 2$ sind nicht beliebig, sondern genau die Z$_3$-Charaktere.
\end{itemize}

\subsection{Was die Algebra für Doc 170 nicht leistet}

\begin{important}[Grenze der Anschlussfähigkeit]
Die Z$_3^3$-Skalenleiter erzwingt \emph{Diskretheit}, aber nicht die konkrete \emph{Stufenzuordnung} aus Doc~170. Die Algebra sagt nicht: \glqq Stufe~4 ist Bewusstsein\grqq{}. Welche $n$ welcher biologischen Stufe entspricht, bleibt eine phänomenologische Frage und ist in Doc~170 als Arbeitshypothese markiert.
\end{important}

Die Drittelpotenzen erzeugen unendlich viele Schichten, nicht nur fünf. Doc~170 wählt fünf, weil die biologische Phänomenologie diese Stufen nahelegt — nicht, weil die Algebra sie ausweist.

\section{Methodische Selbstkontrolle}\label{sec:methodik}

In Doc~206 wurden zwei Selbstkontroll-Abschnitte eingeführt: ein Numerologie-Filter (\S12) und ein methodischer Status-Abschnitt zu kosmologischen Daten (\S11). Dieselbe Disziplin gilt für die Bewusstseinsfrage.

\subsection{Numerologie-Filter}

Wir prüfen, ob unsere Brücken den Status eines theoretischen Befundes haben oder ob sie zufällig wirken könnten.

\begin{itemize}
\item \textbf{Brücke~A (Tekermen $\equiv$ $\xi > 0$):} Algebraisch identisch, nicht numerologisch. Status: theoretisch.
\item \textbf{Brücke~B (Phillips IGL $\equiv$ $A^2 = A^\top$):} Algebraisch identisch (Doc~206 \S7). Status: theoretisch.
\item \textbf{Brücke~C (Doc~170-Stufen $\equiv$ $(1-\xi)^{n/3}$):} \emph{Nur teilweise} theoretisch. Diskretheit und Hierarchie sind algebraisch verankert. Die konkrete Stufenzuordnung ist \emph{nicht} algebraisch gerechtfertigt und bleibt phänomenologisch.
\end{itemize}

\subsection{Was wir nicht behaupten}

Wir vermeiden vier typische Fehlschlüsse:

\begin{enumerate}
\item \textbf{Numerologische Bewusstseins-Identifikation.} Aussagen der Form \glqq Bewusstsein entspricht der Schicht $n = 8$, weil $(1-\xi)^{8/3}$ eine bestimmte Größenordnung hat\grqq{} sind nicht zulässig. Es gibt keinen geometrischen Hinweis auf einen ausgezeichneten $n$-Wert.
\item \textbf{Zirkuläre Adlam-Auflösung.} Wir behaupten nicht, dass $\xi > 0$ das Adlam-Problem \glqq löst\grqq{}. Wir sagen nur, dass es die strukturelle Vorbedingung erfüllt, die Adlam vermisst. Was über dieser Vorbedingung steht, ist offen.
\item \textbf{IGL-zu-Bewusstsein-Sprung.} Aus $A^2 = A^\top$ folgt nicht Bewusstsein. Self-Dual ist notwendig, nicht hinreichend.
\item \textbf{Stufen-Festlegung.} Doc~170 ordnet fünf Stufen zu. Doc~206 erzeugt unendlich viele. Wir lassen offen, welche Stufen \glqq besetzt\grqq{} sind und ob die Doc~170-Zuordnung exakt richtig ist.
\end{enumerate}

\subsection{Methodischer Status der Bewusstseinsfrage}

Doc~206 ist ein bewiesenes Theorem mit reproduzierbaren Skripten. Doc~100 und Doc~170 sind explizit als Anschluss-Schriften und Arbeitshypothesen formuliert. Doc~207 (dieses Dokument) ist eine Synthese-Schrift; es liefert keine neuen Beweise, sondern stellt strukturelle Identifikationen her.

\begin{interpretation}
Die Bewusstseinsfrage ist innerhalb der FFGFT-Reihe nicht abgeschlossen. Doc~206 macht die Sprache der Diskussion algebraisch konsistent, aber die ontologische Frage \glqq Was ist Bewusstsein?\grqq{} bleibt offen. Wer aus FFGFT eine Bewusstseinstheorie ableiten will, muss zwischen strukturellen Aussagen (algebraisch belegt) und phänomenologischen Aussagen (Arbeitshypothese) klar trennen.
\end{interpretation}

\section{Empirische Anschlussfähigkeit: Vorhersagen unter expliziten Zusatzannahmen}\label{sec:vorhersagen}

Die Algebra in Doc~206 ist eine reine Strukturaussage; sie liefert keine direkten Vorhersagen über messbare Bewusstseinsmaße. Wenn man jedoch die \emph{Diskretheits-Hypothese} aus Dokument~170 (Bewusstseinsstufen sind topologisch erzwungen, nicht graduell erreichbar) als Zusatzannahme hinzunimmt, ergeben sich vier qualitative empirische Vorhersagen. Wir formulieren sie in der Form, die Doc~206 \S12 zulässt: \emph{Strukturmerkmale ohne Zahlenwerte}.

\subsection{Was folgt aus FFGFT, was nicht}\label{sec:vorhersagen_folgt}

\begin{important}[Trennung von Strukturaussagen und Zahlenwerten]
Aus FFGFT \emph{folgt}: Wenn die Diskretheits-Hypothese stimmt, sollten Bewusstseinsmaße qualitativ diskontinuierliche Strukturen zeigen (Bimodalität, Sprünge, Kipp-Punkte, taxonomische Schwellen).

Aus FFGFT \emph{folgt nicht}: konkrete Zahlenwerte für Cutoffs, kritische Dosen, Schwellenzeitpunkte oder welche Spezies welche Schwelle überschreiten. Diese hängen von neurobiologischer Phänomenologie ab, nicht von $\xi$.
\end{important}

\subsection{Vorhersage A: Bimodalität neuronaler Bewusstseinsmaße}\label{sec:vorhersageA}

Wenn die Diskretheits-Hypothese stimmt, sollte ein Bewusstseinsmaß wie der Perturbational Complexity Index (PCI) keine kontinuierliche Verteilung zeigen, sondern eine \emph{bimodale} Trennung in zwei Cluster mit einer Lücke dazwischen.

\begin{itemize}
\item \textbf{Folgt:} Bimodalität als qualitatives Strukturmerkmal.
\item \textbf{Folgt nicht:} Konkreter Cutoff (z.\,B.\ PCI~$\approx 0{,}31$, wie er klinisch verwendet wird) oder Lückengröße in $\xi$-Einheiten. Aussagen der Form \glqq Lücke ist proportional zu $\xi$\grqq{} sind ohne algebraische Stütze.
\end{itemize}

\textbf{Empirischer Stand:} Casali et al.~(2013)~\cite{Casali2013} haben einen klinischen PCI-Cutoff etabliert. Ob die Verteilung tatsächlich bimodal ist (echte Lücke) oder unimodal mit pragmatischem Cutoff, ist eine offene empirische Frage.

\subsection{Vorhersage B: Entwicklungssprung statt Gradient}\label{sec:vorhersageB}

Bewusstseinsentwicklung beim Säugling sollte einen Sprung zeigen, keinen kontinuierlichen Gradienten.

\begin{itemize}
\item \textbf{Folgt:} Sprung-Charakter als Strukturmerkmal.
\item \textbf{Folgt nicht:} Zeitpunkt des Sprungs. Dieser hängt von neuraler Reifung ab (z.\,B.\ Myelinisierung rekurrenter Bahnen), nicht von $\xi$.
\end{itemize}

\subsection{Vorhersage C: Anästhesie-Kippen statt Graduierung}\label{sec:vorhersageC}

Bei langsamer Dosissteigerung eines Anästhetikums sollte der Bewusstseinsverlust an einem Kipp-Punkt erfolgen, nicht graduell.

\begin{itemize}
\item \textbf{Folgt:} Kipp-Charakter.
\item \textbf{Folgt nicht:} Kritische Dosis, Substanz-Spezifität, Dosis-Wirkungs-Zeitkonstanten.
\end{itemize}

\subsection{Vorhersage D: Spezies-Klassifikation statt Skalierung}\label{sec:vorhersageD}

Bewusstsein sollte nicht eine Funktion der Komplexität allein sein. Komplexitätszuwachs (mehr Neuronen, mehr Verbindungen) ohne strukturelle Z$_3$-Selbstinversion sollte keine Bewusstseinsschwelle überschreiten lassen.

\begin{itemize}
\item \textbf{Folgt:} Existenz einer strukturellen Schwelle, die durch Skalierung allein nicht überwindbar ist.
\item \textbf{Folgt nicht:} Welche Taxa über/unter der Schwelle stehen. Diese Frage ist neurobiologisch zu beantworten, nicht aus $\xi$ ableitbar.
\end{itemize}

\textbf{Konsequenz für KI:} Eine reine Skalierung von Transformer-Architekturen erreicht Bewusstsein per Vorhersage~D nicht. Notwendig wäre eine Z$_3$-äquivalente Selbstinversion mit metabolischer Verkörperung im Sinne von Doc~170 \S2 (Verwundbarkeit, Metabolismus, $T = 1/m$). Dies ist eine strukturelle Aussage, keine zeitliche Prognose über künftige KI-Entwicklungen.

\subsection{Was wir nicht vorhersagen (Filter analog Doc 206 \S12)}\label{sec:vorhersagen_filter}

Wir grenzen explizit ab, was \emph{nicht} aus FFGFT abgeleitet werden kann, auch wenn es naheliegend erscheint:

\begin{enumerate}
\item \textbf{Keine quantitative PCI-Schwelle aus $\xi$.} Aussagen der Form \glqq $C_{\text{krit}} = \kappa \cdot \xi$ mit $\kappa = \mathcal{O}(1)$\grqq{} sind Numerologie. Ein freier Faktor $\kappa$ erlaubt jeden gewünschten Wert.
\item \textbf{Keine Ableitung des EEG-Spektralexponenten.} Vorschläge wie $\beta = 2/3 + \xi/(2\pi)$ enthalten Faktoren ($1/(2\pi)$), die in der FFGFT-Algebra nicht begründet sind. Der Effekt wäre zudem kleiner als das EEG-Rauschen und damit praktisch unfalsifizierbar.
\item \textbf{Keine Identitätsthese.} Aussagen der Form \glqq Bewusstsein $\equiv$ persistente rekursive Kopplung mit $T = 1/m$-Modulation\grqq{} lösen das Hard~Problem nicht; sie definieren es um. Self-Dual, Nichtabschluss und n/3-Hierarchie sind notwendige strukturelle Bedingungen --- keine hinreichenden.
\item \textbf{Keine Skalenleiter $(1-\xi)^{k-1}$.} Die einzige in Doc~206 algebraisch gestützte Skalenleiter ist $(1-\xi)^{n/3}$. Sie unterscheidet nicht \glqq Bewusstseinsstufen\grqq{} gegen \glqq Nicht-Bewusstseinsstufen\grqq{}.
\item \textbf{Keine Z$_3$-Identifikation mit drei anatomischen Hirnarealen.} Doc~206 hat ein Z$_3$-Dreieck mit drei Knoten in einem abstrakten Skalenraum. Eine Identifikation mit konkreter Anatomie (V1 $\to$ V4 $\to$ präfrontal) ist ein Kategorienwechsel ohne algebraische Stütze.
\end{enumerate}

\subsection{Falsifizierbarkeit}\label{sec:vorhersagen_falsifikation}

Die vier qualitativen Vorhersagen (A--D) sind falsifizierbar:

\begin{itemize}
\item Wenn PCI-Verteilungen empirisch eindeutig kontinuierlich sind, fällt Vorhersage~A.
\item Wenn Säuglings-Bewusstseinsmarker einen sauber kontinuierlichen Gradient zeigen, fällt Vorhersage~B.
\item Wenn Anästhesie-Effekte exakt graduell verlaufen, fällt Vorhersage~C.
\item Wenn Komplexität allein (z.\,B.\ in skalierter KI) Bewusstsein erzeugt, fällt Vorhersage~D.
\end{itemize}

\begin{interpretation}
Eine Falsifikation der qualitativen Vorhersagen falsifiziert die \emph{Diskretheits-Hypothese aus Doc~170}, nicht das algebraische Skelett aus Doc~206. Die Z$_3$-Algebra ist von der Bewusstseinsfrage unabhängig richtig oder falsch.
\end{interpretation}

\section{Verhältnis zur IPI-Mailingliste-Diskussion}\label{sec:ipi}

In den Wochen vor diesem Dokument wurden in der IPI-Mailingliste mehrere zentrale Begriffe vorgeschlagen, die sich nun direkt in das Z$_3$-Skelett aus Doc~206 einordnen lassen.

\begin{itemize}
\item \textbf{Tekermens \glqq open fraction\grqq{}} (Mai 2026): Die Idee, dass Bewusstsein ein \emph{nicht-binäres Spektrum} sei, korrespondiert mit der Z$_3^3$-Skalenleiter. \glqq Nicht-binär\grqq{} bedeutet algebraisch: nicht nur $0$ und $1$, sondern $(1-\xi)^{n/3}$ für alle $n$.
\item \textbf{Phillips' IGL und Self-Dual~C-25} (Mai 2026): Algebraisch identifiziert mit $A^2 = A^\top$ in Doc~206 \S7.
\item \textbf{Austin's $D = 3 + \varepsilon$} (Mai 2026): In Doc~206 \S5 als Spiegelbild von $D_f = 3 - \xi$ identifiziert ($\varepsilon = -\xi$). Auf die Bewusstseinsfrage hat das keine direkte Auswirkung, ist aber methodisch relevant: dieselbe geometrische Struktur kann mit unterschiedlichen Vorzeichen-Konventionen formuliert werden.
\end{itemize}

\subsection{Was die IPI-Diskussion und Doc 207 leisten}

Die Diskussion in der IPI-Mailingliste hat in den Wochen vor diesem Dokument die Sprachebene erreicht, auf der algebraische Identifikationen möglich werden. Doc~207 ist nicht der Anfang dieser Diskussion, sondern eine Konsolidierung.

Die Sprache, in der jetzt über Bewusstsein gesprochen wird (offene Rekursion, Self-Dual, nicht-binäres Spektrum), ist mit der Z$_3$-Geometrie konsistent. Das ist ein methodischer Fortschritt — kein inhaltlicher.

\section{Bilanz}\label{sec:bilanz}

\begin{conclusionBox}[Bilanz: Was Doc 207 herstellt]
Drei strukturelle Brücken zwischen Doc~206 und der bestehenden Bewusstseinsdiskussion sind algebraisch belegt:

\begin{enumerate}
\item \textbf{Konstitutive Offenheit} (Tekermen) $\equiv$ $\xi > 0$ (Nichtabschluss-Theorem aus Doc~193, algebraisch in Doc~206 \S6).
\item \textbf{Information Grounding Law} (Phillips) $\equiv$ $A^2 = A^\top$ (Z$_3$-Selbstinversion in Doc~206 \S7).
\item \textbf{Diskrete Komplexitätsstufen} (Doc~170) algebraisch konsistent mit der Z$_3^3$-Skalenleiter $(1-\xi)^{n/3}$ — Diskretheit erzwungen, Stufenzuordnung phänomenologisch.
\end{enumerate}

Adlams No-Go-Theorem und FFGFTs Nichtabschluss-Theorem sind zwei Seiten derselben strukturellen Aussage.
\end{conclusionBox}

\subsection{Was offen bleibt}

\begin{itemize}
\item Die ontologische Bewusstseinsfrage (\glqq Hard Problem\grqq{}) ist durch keine der drei Brücken berührt.
\item Die konkrete Zuordnung der Doc~170-Stufen zu $n$-Werten ist nicht algebraisch gerechtfertigt.
\item Hinreichende Bedingungen für Bewusstsein bleiben offen; Doc~206 liefert nur notwendige strukturelle Bedingungen für die Sprache der Diskussion.
\item Die Zelle als \glqq erste Bewusstseinsstufe\grqq{} (Doc~170) ist eine Arbeitshypothese, nicht ein Resultat aus Doc~206.
\end{itemize}

\subsection{Verhältnis zum Korpus}

\begin{center}
\begin{tabular}{c|p{8cm}|l}
\textbf{Dokument} & \textbf{Beitrag zur Bewusstseinsfrage} & \textbf{Status} \\
\hline
100 & Bewusstsein als fraktale Inkohärenz, Anschluss Adlam et al. & Anschluss-Schrift \\
170 & Schwellenleiter, $K_{\text{frak}}$ als Schwellenwertoperator & Arbeitshypothese \\
193 & Nichtabschluss-Theorem (algebraisch in 206 verankert) & Bewiesen \\
206 & Z$_3$-Skelett, sechs Brücken, $D_f = 3-\xi$ & Bewiesen \\
\textbf{207} & \textbf{Strukturelle Brücken zwischen 206 und 100/170} & \textbf{Synthese} \\
\end{tabular}
\end{center}

Doc~207 macht keine neue Bewusstseinsaussage. Es zeigt, dass die bestehenden Aussagen aus Doc~100, Doc~170 und der IPI-Diskussion in einem konsistenten algebraischen Rahmen stehen. Die Bewusstseinsfrage selbst bleibt eine offene Forschungsfrage. Doc~206 bietet die Sprache, in der weitere Arbeit präzise formulierbar wird; es bietet keine Antwort auf das, was Bewusstsein ist.

\begin{philosophical}[Schluss]
Die Frage, was Bewusstsein ist, wird durch eine algebraische Brücke nicht beantwortet. Aber die Frage, in welcher Sprache man danach fragen kann, ohne in numerologische oder rein metaphorische Sackgassen zu geraten, hat durch Doc~206 eine schärfere Form. Mehr ist nicht gewonnen. Weniger ist nicht gewollt.
\end{philosophical}

\section*{Danksagung}

Die hier formulierten Brücken stützen sich auf die Diskussionen in der \emph{Information Physics Institute}-Mailingliste und auf die Arbeiten von Onur Tekermen, Phillips Paul und Peter Austin. Die phänomenologische Linie verdankt sich der Auseinandersetzung mit dem No-Go-Theorem von Adlam, McQueen und Waegell~\cite{Adlam2025}.

\begin{thebibliography}{9}
\bibitem{Adlam2025} E.\,C.~Adlam, K.\,J.~McQueen, M.~Waegell. \emph{Agency cannot be a purely quantum phenomenon.} arXiv:2510.13247 (2025).
\bibitem{Casali2013} A.\,G.~Casali, O.~Gosseries, M.~Rosanova, M.~Boly, S.~Sarasso, K.\,R.~Casali, S.~Casarotto, M.-A.~Bruno, S.~Laureys, G.~Tononi, M.~Massimini. \emph{A theoretically based index of consciousness independent of sensory processing and behavior.} \emph{Sci.\ Transl.\ Med.}\ \textbf{5}, 198ra105 (2013).
\end{thebibliography}

\section*{Lizenz und Verfügbarkeit}

\textbf{GitHub:} \href{https://github.com/jpascher/T0-Time-Mass-Duality}{\texttt{github.com/jpascher/T0-Time-Mass-Duality}}\\
\textbf{Zenodo (Vorgänger v1.0.13):} \href{https://zenodo.org/records/20041529}{\texttt{zenodo.org/records/20041529}} (DOI: \href{https://doi.org/10.5281/zenodo.20041529}{10.5281/zenodo.20041529})\\
\textbf{ORCID:} \href{https://orcid.org/0009-0000-6518-4064}{0009-0000-6518-4064}\\
\textbf{Lizenz:} CC-BY-SA 4.0


\end{document}
